% Model-description, calibration and consistency audit of the trend-growth branch.
% Compile: latexmk -pdf model_audit_2026-10-01.tex   (from code/reports/)
% Table bodies and the dead-code inventory are snapshots in audit_2026-10-01/.
\documentclass[11pt,a4paper]{article}
\usepackage[margin=2.1cm,top=2cm,bottom=2cm]{geometry}
\usepackage{amsmath,amssymb,booktabs,longtable,array,calc,siunitx,textcomp,enumitem,xcolor}
\usepackage[labelfont=bf,font=small,skip=4pt]{caption}
\usepackage[hidelinks]{hyperref}
\usepackage{underscore}
\sisetup{table-align-text-post=false}
\setlength{\parskip}{3pt plus 1pt}
\setlength{\parindent}{0pt}
\setlist{nosep,leftmargin=*}
\providecommand{\tightlist}{\setlength{\itemsep}{0pt}\setlength{\parskip}{0pt}}
\newcommand{\np}{\ensuremath{\text{--}}}
\newcommand{\tabdir}{audit_2026-10-01}
\newcommand{\code}[1]{\texttt{#1}}
\newcommand{\stale}{\textcolor{red!70!black}{\textsc{stale}}}
\newcommand{\verified}{\textsc{verified}}
\newcommand{\ok}{\textsc{ok}}
\renewcommand{\arraystretch}{1.08}

\title{\vspace{-1.4cm}Model description, calibration and consistency audit\\[2pt]
\large Branch \code{trend-growth} at \code{3784ae5} (working tree \code{dce6352}), 2026-10-01}
\author{}
\date{}

\begin{document}
\maketitle
\vspace{-1.2cm}

\section*{Object and method}

\textbf{Object.} The code under \code{code/} at commit \code{3784ae5} (2026-10-01 22:32); the working tree is \code{dce6352}, which differs from it only by the audit plan file. Configuration \code{calibration_input_GR.json} (modified 22:25), the data files \code{data/demography_GR.npz}, \code{data/survival_GR.npz}, \code{data/europop2023_GR.npz}, and \code{data/pension_floor_GR.json}. The baseline is the demographic transition from the measured 2023 cross-section with no policy change. The fiscal layer is in scope as equations; no G or $I_g$ experiment result exists under trend growth and none is audited.

\textbf{Method.} Six context-free agents read the source, the configuration and the data files in parallel (22:53--23:15), each on one block of questions: household problem; initial condition and cohorts; aggregation, transition and terminal state; calibration procedure; status of the six items of \code{code/docs/OPEN_ISSUES_2026-07-30.md}; dead code. They were forbidden the draft, the project notes, the earlier code report, the model sections of the README, the assistant's memory and the git history; \code{code/CLAUDE.md} was moved out of the tree before launch and restored after the last agent returned (\code{git status} clean apart from the pre-existing \code{docs} submodule pointer). Every Critical and Major claim below was re-read in the source by the author of this report; the dead-code searches for the unreachable items were repeated. Local numerical checks ran on an 8-core Apple CPU (JAX 0.8.0, float64); the GPU runs of the plan (C1, C1b, C2) were not run, since no instance was available in the session; a reduced-scale substitute for C1 ran overnight on the CPU. Appendix~C lists every check with its outcome.

\textbf{Reading guide.} Part~I states the model as the code implements it, in the notation of the draft's model section where the objects coincide, without parameter values and without solution details. Part~II states the solution method, the calibration and the provenance of every recorded number. Part~III ranks the findings, reports the status of the six July items, summarises the dead-code inventory, lists what could not be verified, records the comments received on the first draft, and lists the changes made after the audit. The appendices give the code-to-equation correspondence, the dead-code inventory, and the checks.

%====================================================================
\section{The model economy as implemented}

Time $t = 0, 1, \dots$ indexes periods, period $t$ being calendar year $2023 + t$, over a horizon $T_{tr}$. Ages are $j = 1, \dots, J$ with real age $24 + j$; working ages $\mathcal{W} = \{1, \dots, J_R\}$ and retired ages $\mathcal{R} = \{J_R + 1, \dots, J\}$. Education $e \in \{\text{low}, \text{medium}, \text{high}\}$ with fixed population shares $s_e$. Labour productivity grows at the exogenous rate $g$; every quantity below is detrended by $Z_t = (1+g)^t$, and aggregates are per capita, meaning per person aged $25$ to $24 + J$ alive in the period. The hats of the draft are dropped.

\subsection{Demographics}

A cohort enters at age 1. The cohort entering in year $y$ has size $N^0_y$ and faces at age $j$ the survival probability $\pi_j(y)$ read from the period life table of the calendar year in which it reaches that age: the observed tables through 2023, the EUROPOP2023 baseline mortality for 2024--2100, and the 2100 table for every later year. Cumulative survival to age $j$ is $S_j(y) = \prod_{i<j} \pi_i(y)$. The sizes of the entering cohorts are set as follows. For years up to 2023 they are recovered from the measured 2023 population: the number of people of each age in 2023 divided by their cumulative survival, so that the model's 2023 population equals the measured one age by age. For 2024--2100 the entering cohort is the number that makes the model's population aged $25$ to $24 + J$ equal to the EUROPOP2023 projection, which books net migration at the entry age. After 2100 the growth rate of the entering cohort is brought linearly to zero between 2100 and 2120 and held at zero afterwards. The age distribution of the population is then constant from 2179, one lifetime later.

The population alive in period $t$ is $N_t = \sum_j N^0_{t-j+1} S_j(t-j+1)$, its growth rate $n_t = N_{t+1}/N_t - 1$, and the growth factor of every detrended per-capita stock is
\[
\Gamma_t = (1+g)(1+n_t), \qquad \Gamma_t = \Gamma_T \equiv (1+g)(1+n_\infty) \ \text{for } t \ge 156 \ (\text{year } 2179).
\]

A per-capita aggregate is the population-share-weighted sum of cohort averages: with $\tilde{x}_{e,j}(t)$ the average of a variable over the surviving members of the cohort aged $j$ in period $t$ and education $e$,
\[
X_t = \sum_e s_e \sum_{j=1}^{J} \frac{N^0_{t-j+1}\, S_j(t-j+1)}{N_t}\; \tilde{x}_{e,j}(t).
\]
Because a cohort's share of the population one period later is its current share divided by $(1+n_t)$, the household's detrending by $(1+g)$ and the stocks' by $\Gamma_t$ are consistent.

\subsection{Assets and returns}

Households hold one asset $a \ge 0$ that earns the world return $r$ whatever it finances, taxed at $\tau^k$ on the interest. Domestic private capital $K_t$ is chosen by firms (\S\ref{sec:tech}); sovereign debt $B_t$ is held abroad and serviced at the sovereign rate $r_B$. The household sector's foreign assets are $A_t - K_t$, and the economy's net foreign asset position nets the government's external debt out of them,
\[
\mathrm{NFA}_t = A_t - K_t - B_t, \qquad A_t = \text{per-capita household wealth.}
\]
There is no portfolio choice and no risk premium; $r$ and $r_B$ are two exogenous prices. The assets of households who die are accidental bequests. The model has a bequest tax $\tau^{beq}$ and a lump-sum transfer of the remainder to the entering cohort; at the audited commit both are switched off, so accidental bequests leave the economy (finding M5). Since 2026-10-02 the configuration taxes them away in full, $\tau^{beq} = 1$ (\S\ref{sec:postaudit}).

\subsection{Prices and policy}

The interest rate $r_t$ is exogenous and the wage $w_t$ is the firm's marginal product of labour given the public capital stock (\S\ref{sec:tech}). The policy instruments are:
\begin{itemize}
\item four proportional tax rates: $\tau^c$ on consumption, $\tau^l$ on labour income net of social contributions, on unemployment benefits and on pensions, $\tau^p$ (social contributions) on wage income only, and $\tau^k$ on interest income;
\item a pay-as-you-go pension paid from the general budget: a replacement rate $\rho$ applied to a pension base built from the retiree's last income state and the career average (equation \eqref{eq:pens}), subject to a minimum pension $b_{\min}$; the benefit is constant in detrended units, so in levels it grows with productivity;
\item unemployment insurance replacing a fraction $\rho^{ui}$ of the earnings the household's previous income state would have delivered, for the first period of a spell;
\item public coverage of a fraction $\kappa$ of each person's age-dependent medical expenditure $m(j)$, the rest paid out of pocket;
\item government consumption $G_t$, defence spending $D_t$ and other net spending $O_t$, each a fixed share of output; $O_t$ is the balancing item, the net of expenditure and revenue items the model does not represent, set once so that the base-year primary balance equals the data and held as a share of output thereafter;
\item public investment $I^g_t$, set in the baseline at the level that holds detrended per-capita public capital constant;
\item sovereign debt, which absorbs the primary balance at rate $r_B$. In the tax-financed experiments the labour income tax rate is raised by a constant amount from $t = 0$, chosen so that the economy reaches a target debt ratio, or a target net-foreign-asset ratio, at the end of the horizon. No other instrument moves.
\end{itemize}

\subsection{The household problem}

\textbf{State and shocks.} $(e, \alpha, j, z, z_{last}, a)$. The permanent productivity effect $\alpha \sim N(0, \sigma^{2}_{\alpha,e})$ is drawn at entry. The income state $z \in \{0, z_1, \dots, z_{n-1}\}$: $z = 0$ is unemployment; the employed states follow an autoregressive process in logs with persistence $\rho_z$ and an education-specific mean $\mu_e$ and innovation variance $\sigma^2_{\varepsilon,e}$. From unemployment a job is found with probability $f$, the new state being drawn with equal probability from the employed states; from employment a separation occurs with probability $s_e$, set so that the stationary unemployment rate equals the education-specific rate $u_e$. $z_{last}$ is the previous period's state, whatever it was, and is frozen at retirement. Health has one state; medical expenditure $m(j) = m\,\tilde{m}_j$ is an age profile scaled by a constant. At entry $a = 0$, $z$ is drawn from the stationary distribution of the income process, and $z_{last} = z$.

\textbf{Preferences.} $u(c, \ell) = \log c - \nu\,\ell^{1+\varphi}/(1+\varphi)$. Hours $\ell \in [0, 1]$ are chosen by the employed; the unemployed and the retired supply $\ell = 0$. Discounting is $\beta\,\pi_j$; death leaves no bequest motive. With log consumption the detrended problem needs no growth adjustment of $\beta$.

\textbf{Budget constraints}, in detrended units, with $\kappa_j$ the age-productivity profile and $y^L = w\,\kappa_j\,z\,e^{\alpha}$ the wage per unit of hours:
\begin{align}
\text{working:}\quad & (1+\tau^c)\,c + (1+g)\,a' = \bigl(1 + (1-\tau^k) r\bigr) a + (1-\tau^l)\bigl[(1-\tau^p)\, y^L \ell + T^{UI}\bigr] - (1-\kappa)\, m(j), \label{eq:bcw}\\
\text{retired:}\quad & (1+\tau^c)\,c + (1+g)\,a' = \bigl(1 + (1-\tau^k) r\bigr) a + (1-\tau^l)\,\mathrm{PENS} - (1-\kappa)\, m(j), \label{eq:bcr}\\
& T^{UI} = \mathbf{1}\{z = 0\}\,\rho^{ui}\, w\, \kappa_j\, z_{last}\, e^{\alpha}, \qquad
\mathrm{PENS} = \max\Bigl\{\rho\, w_{J_R}\bigl[\lambda\,\kappa_{J_R} z_{last} + (1-\lambda)\,\bar{\kappa}\,\bar{z}_e\bigr] e^{\alpha},\; b_{\min}\Bigr\}, \label{eq:pens}
\end{align}
with $a' \ge 0$, $w_{J_R}$ the cohort's wage in its last working year, $\bar{\kappa}$ the mean of the age profile over working ages and $\bar{z}_e$ the mean of the employed income states. The weight $\lambda = (1 - \hat\rho^{J_R})/(J_R(1-\hat\rho))$ is the coefficient of a career average on its terminal value for an autoregressive process with persistence $\hat\rho$ over $J_R$ periods; $\hat\rho$ is a separate constant, not the $\rho_z$ of the income process, and the two are equal in the current calibration. Because $z_{last}$ is the previous state including $z = 0$, UI is paid for one period per spell, an entrant drawn unemployed receives none, and a household unemployed at $J_R$ retires on the $(1-\lambda)$ term alone. There is no means-tested consumption floor and no bequest is received.

\textbf{Labour supply.} For an employed household and each candidate $a'$, hours satisfy
\[
\nu\,\ell^{\varphi}\,(1+\tau^c) = \frac{1}{c}\,\mathrm{MW}, \qquad \mathrm{MW} = w\,\kappa_j\, z\, e^{\alpha}\,(1-\tau^p)(1-\tau^l),
\]
with $c$ the consumption implied by \eqref{eq:bcw} at those hours and that $a'$: the marginal disutility of an hour equals the marginal utility of the after-tax wage it buys. The solution is restricted to $[0, 1]$ (\S\ref{sec:solution}).

\textbf{Bellman equations.} Working age:
\[
V_j(z, z_{last}, a) = \max_{a', \ell}\; u(c, \ell) + \beta\,\pi_j\, \sum_{z'} P_z(z' \mid z)\, V_{j+1}(z', z, a'), \qquad \text{s.t. } \eqref{eq:bcw}.
\]
Retired age, as implemented in both backends:
\[
V^R_j(z_{last}, a) = \max_{a'}\; u(c, 0) + \beta\,\pi_j\, V^R_{j+1}(0, a'), \qquad \text{s.t. } \eqref{eq:bcr},
\]
that is, the continuation value is evaluated at $z_{last} = 0$ for every retiree, while the pension in \eqref{eq:bcr} and in the simulation depends on the retiree's own $z_{last}$ (finding C1). At the terminal age $a' = 0$. A household whose resources are non-positive, which happens at $a = 0$ when unemployed with $z_{last} = 0$, since it has no benefit and owes its out-of-pocket medical cost, is handled by a default rule described in \S\ref{sec:solution} (finding M1).

\textbf{Death.} After consumption and taxes, the household dies with probability $1 - \pi_j$; its beginning-of-period assets are recorded as an accidental bequest.

\subsection{Technology}\label{sec:tech}

\[
Y_t = A\,(K^g_t)^{\eta_g}\, K_t^{\alpha}\, L_t^{1-\alpha}, \qquad
L_t = \sum_e s_e \sum_{j \in \mathcal{W}} \frac{N^0_{t-j+1} S_j(t-j+1)}{N_t}\; \widetilde{\bigl(\kappa_j\, z\, e^{\alpha}\, \ell\bigr)}_{e,j}(t),
\]
labour input being the per-capita sum of hours in efficiency units. In the implementation $L_t$ is obtained by dividing wage income plus unemployment benefits by the wage, which equals the expression above plus $\mathrm{UI}_t/w_t$ (finding M2). Given $r$, the firm's conditions pin the capital--labour ratio, the wage and domestic capital,
\[
\frac{K_t}{L_t} = \left(\frac{r + \delta}{\alpha A (K^g_t)^{\eta_g}}\right)^{1/(\alpha - 1)}, \qquad
w_t = (1-\alpha)\, A\,(K^g_t)^{\eta_g} \left(\frac{K_t}{L_t}\right)^{\alpha}, \qquad K_t = \frac{K_t}{L_t}\, L_t,
\]
so $K_t/Y_t = \alpha/(r+\delta)$ in every period. Public capital follows
\[
K^g_{t+1} = \frac{(1-\delta_g)\, K^g_t + I^g_t}{\Gamma_t},
\]
and the baseline investment level $I^g_t = (\delta_g + \Gamma_t - 1)K^g_0$ keeps it constant.

\subsection{Government budget constraints}

With $\Sigma$ the per-capita aggregation of \S1.1,
\begin{align*}
\mathrm{Rev}^c_t &= \tau^c\, C_t, &
\mathrm{Rev}^l_t &= \tau^l\bigl[(1-\tau^p)\mathcal{B}^{lab}_t + \mathrm{UI}_t + \mathrm{PENS}^{out}_t\bigr], &
\mathrm{Rev}^p_t &= \tau^p\,\mathcal{B}^{lab}_t, &
\mathrm{Rev}^k_t &= \tau^k\, r\, A_t,\\
\mathrm{UI}_t &= \Sigma\, T^{UI}, &
\mathrm{PENS}^{out}_t &= \Sigma\, \mathrm{PENS}, &
\mathrm{HSub}_t &= \kappa\,\Sigma\, m(j), &
\mathcal{B}^{lab}_t &= \Sigma\, y^L \ell,
\end{align*}
\begin{align}
\mathrm{PD}_t &= \mathrm{UI}_t + \mathrm{PENS}^{out}_t + \mathrm{HSub}_t + G_t + I^g_t + D_t + O_t - \bigl(\mathrm{Rev}^c_t + \mathrm{Rev}^l_t + \mathrm{Rev}^p_t + \mathrm{Rev}^k_t\bigr), \label{eq:pd}\\
B_{t+1} &= \frac{(1+r_B)\, B_t + \mathrm{PD}_t}{\Gamma_t}, \qquad B_0 = \bar{b}\; Y_0, \label{eq:debt}
\end{align}
with $\bar{b}$ the base-year debt ratio (in the implementation $Y_0$ is taken from a separate low-precision run, finding M9). Interest $r_B B_t$ is reported alongside the primary balance. A pension fund $S_{t+1} = [(1+r_B) S_t + \mathrm{Rev}^p_t - \mathrm{PENS}^{out}_t]/\Gamma_t$, $S_0 = 0$, is a memorandum item that feeds nothing. The net foreign asset position evolves as $\mathrm{CA}_t = \Gamma_t \mathrm{NFA}_{t+1} - \mathrm{NFA}_t$ with $\mathrm{NFA}_t = A_t - K_t - B_t$.

\textbf{Financing rules and terminal conditions.} Debt financing leaves $\tau^l$ unchanged and lets $B_t$ follow \eqref{eq:debt}. Tax financing adds a constant $\Delta\tau^l$ from $t = 0$, found by a root finder, so that either $B_{T-1}/Y_{T-1}$ equals a target debt ratio or $\mathrm{NFA}_{T-1}/Y_{T-1}$ equals a target net-foreign-asset ratio at the end of the horizon; the targets are the baseline's own terminal ratios. After the horizon the simulation continues for twenty periods with every path held at its terminal value (finding C2). A rest-point condition $\mathrm{PD}_t/Y_t = (\Gamma_T - 1 - r_B)\, b$ exists in the code but is not used. A terminal-convergence check compares the last two periods of $K$, $K^g$, $L$, $Y$, $C$, $A$, NFA and $S$; $B$ is not among them (finding M4).

\subsection{Equilibrium and the detrended balanced-growth path}

Each cohort solves its lifecycle problem given the sequences of prices, taxes and transfers it faces over its own lifetime, under the mortality of its own birth year. Cohorts already alive in 2023 are assumed to have lived their earlier years at the 2023 prices and policies, starting from zero assets at entry; their wealth in 2023 is therefore the wealth that a lifetime at those prices generates, not a measured distribution. Cohorts entering after the horizon face prices and policies held at their terminal values. Aggregates are the population-weighted sums of \S1.1; the wage is set by the firm, domestic capital by the capital--labour ratio, and the net foreign asset position absorbs the difference between household wealth and domestic capital plus sovereign debt. In the baseline, aggregate changes come from the changing age composition of the population and the sizes of the entering cohorts; cohort-level changes come from each cohort's own survival schedule. With tax rates fixed, the age distribution enters no household's problem.

The terminal state is a balanced growth path: from 2179 the age distribution of the population, the growth factor $\Gamma_T$ and every cohort's survival schedule are constant, so the detrended per-capita aggregates are constant; levels grow at $\Gamma_T - 1$ and per-capita quantities at $g$. On that path the debt recursion \eqref{eq:debt} has root $(1+r_B)/\Gamma_T$: when $r_B > \Gamma_T - 1$ a constant debt ratio requires a primary surplus of $(r_B - (\Gamma_T - 1))\,b$ per period, and nothing in the baseline imposes it (finding M4).

The economy's resource constraint, per capita and detrended, with $M_t = \Sigma\, m(j)$ total medical spending and $\mathrm{Beq}_t$ the accidental bequests of those who die in $t$ (which leave the economy in the baseline), is
\[
\Gamma_t\mathrm{NFA}_{t+1} - \mathrm{NFA}_t = Y_t + r\,\mathrm{NFA}_t + (r - r_B) B_t - C_t - \bigl[\Gamma_t K_{t+1} - (1-\delta)K_t\bigr] - G_t - I^g_t - D_t - O_t - M_t - \mathrm{Beq}_t ,
\]
where $r\,\mathrm{NFA}_t + (r - r_B)B_t$ is net factor income from abroad: households earn $r$ on their foreign assets $A_t - K_t$ while the government pays $r_B$ on $B_t$. The residual that the report generator and the results checker compute omits $M_t$ and $(r - r_B)B_t$, and the implementation's output carries the unemployment-benefit term of finding M2 (finding M3).

%====================================================================
\section{Solution method and calibration}

\subsection{Solution and simulation}\label{sec:solution}

\textbf{Household problem.} Backward induction over age on a fixed asset grid of 100 points, $a_i = 200\,(i/99)^{1.5}$, eight of them below $a = 4$; the savings choice is a search over the grid with consumption as the residual, without interpolation. The employed income states are a four-point Tauchen discretisation of the autoregressive process two standard deviations wide; the permanent effect $\alpha$ is a five-point Gauss--Hermite discretisation of its normal distribution. For each candidate $a'$ the labour first-order condition is solved by a safeguarded Newton iteration on $[0, 1]$ and the solution restricted to that interval, so that where the bound binds the condition holds as an inequality (\S3.7). The unemployed and the retired carry a placeholder $\ell = 1$ in the decision arrays with zero disutility and zero wage income, which is $\ell = 0$ economically; the simulated sample records $\ell = 1$ for the unemployed and $0$ for retirees, and neither value enters any aggregate or moment. A state with non-positive resources is assigned $c = 10^{-10}$, $a' = 0$ and a value with no continuation term (finding M1). A means-tested consumption floor exists as a parameter; the solver refuses a positive value, and the configuration sets zero. Two implementations exist, NumPy and JAX; the production chain uses JAX, which solves the cohorts of an education group together.

\textbf{Simulation and aggregation.} Each cohort and education group is simulated with $n_{sim}$ households from entry. The per-cohort average over the households the cohort had at entry, counting those who have died as zero, equals cumulative survival times the average among survivors, so weighting those averages by entering-cohort sizes and dividing by the living share gives the per-capita aggregate of \S1.1. Different seeds are used per cohort and education group in the transition; in the calibration all three education groups of a cohort share one seed.

\textbf{Two constructions of the 2023 cross-section.} The \emph{single-lifecycle} construction solves one lifecycle problem per education group under one survival schedule and reads the age-$j$ outcome off that same lifetime for every $j$. The \emph{cohort-by-cohort} construction solves one lifecycle problem per 2023 cohort, each on its own survival schedule, and reads each cohort at its own age; it is how the transition builds its $t = 0$ and costs 60 solves per education group where the other costs one. The configuration now selects the cohort-by-cohort construction for the calibration, the normalisation and the closure.

\textbf{Counterfactuals.} Cohorts alive in 2023 keep their baseline decision rules for the years before 2023, so that household wealth at $t = 0$ is the same in every scenario; the construction passes the regression check on both implementations (Appendix~C). The tax-financed experiments find $\Delta\tau^l$ by the Illinois method to a tolerance of $10^{-3}$.

\subsection{Externally set}

Layout as in \code{reports/calibration_report.tex}. The table bodies are snapshots, taken for this audit, of the untracked files in \code{output/calibration_growth/} that \code{reports/fill_report.py} wrote: the parameter, moment, implied and age bodies and the header at 22:23 on 2026-10-01 from the current configuration; the growth body at 14:36 and the figure at 14:46 from a local baseline run whose configuration state is not recorded. Git records none of them. Every model number in them, and every fitted or set value in the configuration, comes from the 12:31 run on an A100 or from the two scripts that followed it, on the single-lifecycle construction, with five fitted parameters and a pension floor of 0.15; the code and the configuration have moved since (Part~III, C3). Numbers that the live pipeline would not reproduce are flagged \stale.

\begin{table}[htbp]\centering\footnotesize
\caption{Parameters as printed by the report generator on 2026-10-01 22:23. Row flags: $n$ is the configuration's terminal rate $n_\infty = 0$, whereas the realised path starts at $n_0 = -0.0048$ (the header of the same run prints $-0.0048$); $K_g/Y$ is 0.703, the 2019 IMF stock rolled to 2023; $b_{\min} = 0.1671$ is the value in force now, while every fitted value in the next block was obtained at 0.15 (\stale\ for the block as a whole); the five SMM rows omit the sixth parameter $\rho^{ui}$ that the configuration lists; $A$ and $O/Y$ were set at the 12:31 $\theta$ on the single-lifecycle construction (\stale).}
\input{\tabdir/params_body.tex}
\end{table}

Values not in the table, from the configuration: $J = 60$, $J_R = 39$, $T_{tr} = 180$; $r = 0.04$, $r_B = 0.019$; $\tau^c = 0.1818$, $\tau^l = 0.10$, $\tau^k = 0.2236$; $\delta_g = 0.04281$; $G/Y = 0.13$, $D/Y = 0.03$, $\bar b = 1.6428$; $s_e = (0.2343, 0.4705, 0.2952)$; $u_e = (0.1645, 0.158, 0.1005)$; $f = 0.5$, with no recorded source in the code or configuration (\S3.7); $\rho_z = 0.95$, $\mu_e = (-0.1529, 0, 0.259)$, $\sigma_{\varepsilon,e} = (0.0538, 0.0858, 0.0750)$, $\sigma_{\alpha,e} = (0.367, 0.259, 0.318)$; $\hat\rho = 0.95$, hence $\lambda = 0.4434$; $\tau^{beq} = 0$ by default; $n_\infty = 0$; $n_{sim} = 2\,000$ per cohort and education group in the transition and in the cohort-by-cohort calibration, $10\,000$ per education group in the single-lifecycle construction. The age profiles $\kappa_j$ and $\tilde m_j$ are 60-entry vectors in the configuration. Derived: $\Gamma_0 = 1.01212$, $\Gamma_T = 1.017$, $K/Y = 3.667$, $w = 2.1054$ at the current $A$; the old-age dependency ratio (65--84 over 25--64) is 0.359 in 2023 and 0.468 on the terminal path; the living share of the ever-entered population is 0.899 in 2023 and 0.970 on the terminal path. The configuration's 60-entry survival vector is the 2020 period table and enters no production computation, since both constructions take cohort-specific schedules from the demography file.

\subsection{Fitted jointly by SMM}

Targets, weights $1/m^2_{data}$, and the parameter each is labelled as identifying (a labelling in the report generator, not a constraint of the optimiser):

\begin{center}\footnotesize
\begin{tabular}{@{}llS[table-format=1.4]S[table-format=5.2]l@{}}
\toprule
Target & Population & {Data} & {Weight} & Labelled parameter \\
\midrule
Average hours & employed working-age households & 0.4100 & 5.95 & $\nu$ \\
$A/Y$ & whole population & 4.0000 & 0.06 & $\beta$ \\
Payroll revenue$/Y$ & $\tau^p \times$ wage bill & 0.1300 & 59.17 & $\tau^p$ \\
Pensions$/Y$ & retirees & 0.1600 & 39.06 & $\rho$ \\
Public health$/Y$ & $\kappa\, m \sum_j \omega_j \tilde m_j / Y$ & 0.0540 & 342.94 & $m$ \\
UI$/Y$ & working-age unemployed & 0.0060 & 27777.78 & $\rho^{ui}$ (never fitted) \\
\bottomrule
\end{tabular}
\end{center}

Means are taken among the living at each age and weighted by the measured 2023 age shares and the education shares. Two targets are identities given $Y$ and UI$/Y$: payroll revenue$/Y = \tau^p(1 - \alpha - \mathrm{UI}/Y)$ because $w L/Y \equiv 1 - \alpha$ and the base excludes UI, and public health$/Y = \kappa\, m \cdot 0.904951/Y$ because medical spending does not depend on the solution; the fitted $\tau^p = 0.197414$ and $m = 0.089724$ reproduce these identities at the 12:31 run's $Y = 0.9954$ and UI$/Y = 0.0115$. The objective is $\sum_i W_i (m_i^{data} - m_i^{model})^2$ minimised by Nelder--Mead in logit-transformed parameters (tolerance $10^{-5}$ in the shell scripts, 500 iterations at most, same seeds at every evaluation); $\theta$ is written to the configuration only on convergence.

\begin{table}[htbp]\centering\footnotesize
\caption{Moments as printed on 2026-10-01 22:23, read from the 12:31 markdown report (\stale): fitted at $A = 1.42370$ before the normalisation moved $A$ to $1.42761$, at which the normalisation log shows $A/Y = 4.0439$ (+1.10\%) and public health$/Y = 0.0537$; UI$/Y$ is filed as not targeted although the configuration targets it, and its deviation cell prints the absolute gap $+0.005$ as 0.01 (relative: +83\%). Five parameters, single-lifecycle construction, $n_{sim} = 10\,000$ per education type, pension floor 0.15, 169 Nelder--Mead iterations, 286 s.}
\input{\tabdir/moments_body.tex}
\end{table}

\textbf{The 12:31 run and what followed.} \code{run_scale_loop.sh} starts the SMM from the stored $\theta$, runs it, then re-normalises $A$ so that base-year output equals one at the fitted $\theta$ (root finder, tolerance $10^{-4}$), and declares the outer loop converged when the residual \emph{before} the update, $Y(\theta_r, A_{r-1}) - 1$, is below $5\times 10^{-3}$; on 2026-10-01 it was $-4.60\times 10^{-3}$ and $A$ moved from 1.42370 to 1.42761 after the test had passed. The pair written to the configuration is therefore $(\theta_r, A_r)$, which no step evaluates jointly. The balancing item was then set at that pair: household-side primary balance $+0.1045$ of output, discretionary spending $0.1986$, target $+0.0195$, hence $O/Y = -0.113636$.

\textbf{Construction used.} The 12:31 fit used the single-lifecycle construction: one lifecycle per education type on the configuration's survival vector, 10\,000 households, every age read off the same households. The configuration now selects the cohort-by-cohort construction for the SMM, the normalisation and the balancing item. The construction is not recorded in any log or in the metadata; the attribution rests on the output levels in the records depending on $n_{sim}$, which the cohort-by-cohort construction ignores. The cohort-by-cohort construction costs 60 solves per education group where the other costs one (Appendix~C, C0).

\subsection{Set outside the SMM}

$A = 1.42761013$ (normalisation $Y = 1$ at the base-year equilibrium, single-lifecycle construction, 2026-10-01 after 12:31; \stale) and $O/Y = -0.113636$ (balancing item, same chain; \stale). The script that sets the balancing item now recomputes $I_g/Y$ from the level the transition spends, $(\delta_g + \Gamma_0 - 1)K^g/Y = 0.03862$, but writes the value formed from the configuration's ratio 0.03862; the two coincide to five decimals today.

\begin{table}[htbp]\centering\footnotesize
\caption{Calibrated parameters against their direct data counterparts. The 22:23 run wrote placeholders (no \code{--implied} flag); the A100 values of 13:41 survive only in \code{step0_report.log}: contribution base$/Y$ 0.6587, contribution rate 0.1974, pension at retirement over last wage 0.5723 against data 0.570, 0.228, 0.763. The first two model values are identities, $1 - \alpha - \mathrm{UI}/Y$ and $\tau^p$; only the replacement row is a simulated statistic.}
\input{\tabdir/implied_body.tex}
\end{table}

\subsection{Demographics}

\begin{table}[htbp]\centering\footnotesize
\caption{Age structure, 2026-10-01 22:23, from the demography file and the transition's population weights; current, no solve involved. Model $t = 0$ reproduces the data by construction of the demography file (Appendix~C, spot checks 1--3).}
\input{\tabdir/age_body.tex}
\end{table}

\begin{table}[htbp]\centering\footnotesize
\caption{Implied growth rates, written 14:36 on 2026-10-01 by a local baseline run at $n_{sim}$ and configuration not recorded (\stale, provenance unknown). The Level and Per-capita columns are $\Gamma_T - 1$ and $g$ by construction; the detrended-trend column over $t = 156, \dots, 179$ is the flatness test. $B$ is absent because that run tracked no debt.}
\input{\tabdir/growth_body.tex}
\end{table}

%====================================================================
\section{Audit}

\subsection{Summary}

The code at the audited commit implements the economy of Part~I. Two defects in the equations as coded are Critical: the retired household's continuation value is evaluated at the wrong pension in both implementations (C1), and the configuration-driven fiscal script cannot complete a baseline because its 200-period horizon overruns the entering-cohort table, which ends in 2210 (C2). The recorded parameter vector, $A$ and balancing item were produced under a pension floor, a parameter count and a construction of the 2023 cross-section that the live pipeline no longer has, so no number in Part~II is reproducible as the code stands (C3). Among the Major items, the labour aggregate includes unemployment benefits, the accidental-bequest circuit is open, the goods-market residual that the results checker fails runs on is not the model's resource constraint, the baseline debt ratio has no rest point under the balancing item as set and the terminal check does not look at debt, non-positive-resource states are handled by a default rule that drops the continuation value, two scripts fail on their own on the sixth parameter, and the report tables carry several internal inconsistencies. Of the six July items, five are unchanged in mechanism and one (the base-year-versus-transition level gap) has had its mechanism replaced for the production configuration; a reduced-scale run on the CPU puts the gap at $+1.2\%$ on the cohort-by-cohort construction. The aggregation, the demographic path, the growth isomorphism of the household problem, the stock recursions and the predetermination of initial wealth across scenarios check out at test scale; the production-scale checks that need a GPU were not run.

\subsection{Findings}

Each item: the fact, its location, the consequence, and how it was verified. ``Both implementations'' means \code{lifecycle_perfect_foresight.py} (NumPy) and \code{lifecycle_jax.py} (JAX).

\subsubsection*{Critical}

\textbf{C1. Retired continuation value indexed at $z_{last} = 0$ (both implementations).} In the retired branch of the Bellman step the continuation is $V_{j+1}(a', z = 0, h', z_{last} = 0)$ --- \code{lifecycle_perfect_foresight.py:993} (\code{self.V[t + 1, i_a_next, 0, i_h_next, 0]}) and \code{lifecycle_jax.py:332} (\code{V_next[:, 0, :, 0]}) --- while the pension in the same period's budget depends on $z_{last}$ with $\lambda = 0.4434$ (\code{:727--733}; JAX \code{:150--154}) and the simulation keeps $z_{last}$ fixed through retirement and pays the pension at the household's own $z_{last}$ (\code{:1270--1271}, \code{:1194--1200}; JAX \code{:790}, \code{:716--720}). Every retiree with $z_{last} > 0$ therefore chooses $a'$ as if its pension base fell next period to $(1-\lambda)\bar{\kappa}\bar{z}_e$: at the mean employed state and $\alpha = 0$ the pension actually paid is 0.457/0.560/0.711 (low/medium/high) against 0.247/0.303/0.385 assumed for next period, a ratio of 0.54 at every education level, with the floor not binding. The working-age continuation at $J_R$ reads the same retired value, so pre-retirement saving is affected too. Consequence: the retiree consumption and asset profiles, aggregate household wealth $A$ (the $A/Y$ target and hence $\beta$), capital-income-tax revenue and the size of accidental bequests are computed from decision rules that solve a different problem from the one simulated. Agent B1's small-model check ($J = 14$, correcting the index to $z_{last}$) changed $a'$ at 86\% of retired states and reversed the slope of retiree consumption. Verified by reading both call sites and the budget and simulation lines; production magnitude not measured. The draft's retired Bellman equation carries $V^R_{j+1}(z_{last}, a')$, so the code disagrees with the draft here.

\textbf{C2. The configuration-driven fiscal script overruns the entering-cohort table.} \code{run_fiscal_figures.py} asks for twenty periods beyond the horizon on every scenario (\code{n_post = 20}, \code{:218} onwards), so \code{run_fiscal_scenario} extends every path and runs the transition for $T_{tr} + 20 = 200$ periods (\code{fiscal_experiments.py:730--736}). The population weights need the entering cohort of year $2023 + t$, and \code{_entrant_weights} raises \code{ValueError} when that year exceeds the table's last year 2210 (\code{olg_transition.py:947--953}; the demography file is built ``to cover T_transition = 180'', \code{build_demography_GR.py:48}). Verified by calling the method on the production economy: $t = 187$ returns, $t = 188$ raises (``entering cohorts cover 1964..2210 but period t=188 needs 2152..2211''); \code{growth_factors} clips the same overrun instead. Consequence: \code{python run_fiscal_figures.py --config calibration_input_GR.json} cannot complete its baseline, so no fiscal number can be produced on the current configuration. No reported number is corrupted: both \code{fiscal_results.json} files in \code{output/} are 60-period runs from July, without the growth and aggregation stamps.

\textbf{C3. No recorded number is reproducible by the live pipeline.} The configuration's $\theta$ (five values), $A = 1.42761013$ and $O/Y = -0.113636$ come from the 12:31 chain with \code{pension_min_floor = 0.15} (every one of the seven calibration reports since 29 September prints 0.15), five parameters and targets, and the single-lifecycle construction; the configuration now has the floor at 0.1671 (\code{build_pension_floor_GR.py}, 22:21), six parameters and targets with UI$/Y$ at weight 27\,777.78, and \code{base_year_cohorts = true}; \code{calibrate.py} and \code{pin_baseline_closure.py} were edited after the run (file times 22:29 and 22:15). In addition $\theta$ was fitted at $A = 1.42370$ and sits beside $A = 1.42761$, at which the model's $A/Y$ is 4.0439 (normalisation log) while the moments table prints 4.0002. Consequence: every model column in Part~II describes a configuration that no longer exists; the direction of the change in $\theta$ is not readable from the code (the test file records that the cohort-by-cohort construction moved hours by +2.6\% and pensions$/Y$ by $-3.1$\% on the earlier production run); at the stored $\theta$ the CPU run at $n_{sim} = 200$ puts hours at $+1.6\%$ and $A/Y$ at $+6.3\%$ of their targets on the live configuration (Appendix~C, one noisy draw). Verified against the configuration, the reports and the logs; the agents' attribution of the construction is an inference from the $n_{sim}$-dependence of the recorded levels.

\subsubsection*{Major}

\textbf{M1. Non-positive resources handled by a default rule with no continuation (both implementations).} When no $a'$ leaves $c > 0$, the solver returns $c = 10^{-10}$, $a' = 0$ and $V = u(c)$ with no $\beta\pi_j\mathbb{E}V_{j+1}$ term (\code{lifecycle_perfect_foresight.py:1020--1024}; \code{lifecycle_jax.py:354--361}). The state $a = 0$, $z = 0$, $z_{last} = 0$ has resources $-(1-\kappa)m(j) \in [-0.010, -0.035]$ because UI is zero there and out-of-pocket medical spending is due. Every entrant starts at $a = 0$ with $z_{last} = z$ and $z$ drawn from the stationary distribution, so the share of entrants in that state at entry equals the education unemployment rate (16.5\%, 15.8\%, 10.1\%); it recurs for any zero-asset worker in the second year of a spell (probability 0.5). Consequences: the medical outlays are recorded although the budget cannot pay them, so the household budget identity fails by the shortfall in those household-periods; $c = 10^{-10}$ enters the consumption aggregate and every consumption-distribution statistic; the value of hitting $a = 0$ while unemployed, against which neighbouring states' precautionary saving is computed, omits the future. Verified by reading the default rule, the entry state and the budget. Share of household-periods not measured.

\textbf{M2. Labour input includes unemployment benefits (July item 1, unchanged).} \code{effective_y_sim = wage_income + ui_sim} (\code{lifecycle_perfect_foresight.py:1232}; \code{lifecycle_jax.py:734}) is the per-age labour mean and is divided by $w$ to form $L$ (\code{olg_transition.py:2173}; \code{calibrate.py:628, :1524}); the social-contribution base is wage income alone (\code{:1236}; JAX \code{:746}). Dividing wage income by the wage gives exactly hours in efficiency units, so the discrepancy from the model's $L_t$ is the benefit term alone: $L$, $K$, $Y$ and $wL$ are overstated by UI$/(wL) \approx 0.011/0.67 \approx 1.7\%$ at the logged UI$/Y$; the same convention holds in calibration and transition, so $\theta$ absorbs it and the levels and the multiplier denominator carry it. Verified.

\textbf{M3. The goods-market residual is not the model's resource constraint.} \code{reports/fill_report.py:302--364} and \code{eval_fiscal_results.py:225--262} compute $C - (Y - I - (G/Y + D/Y + O/Y)Y - I^g - \Delta\mathrm{NFA})$ and an ``open'' form that subtracts $r\,\mathrm{NFA}$. Against the resource constraint in \S1.7, two booked flows are missing: total medical spending $M_t$ (household side \code{:768}, government side \code{olg_transition.py:1824}) and $(r - r_B)B_t$. A third term appears because of M2: the implementation's $Y_t$ and $K_t$ contain the benefit inside $L_t$, while the government also pays it, so the identity on the flows as booked carries an extra $-\mathrm{UI}_t$. On the configuration's own ratios the open form cannot be smaller in magnitude than $M/Y + \mathrm{UI}/Y - (r - r_B)B/Y \approx 0.082 + 0.006 - 0.034 = 0.054$ of output in a debt-financed run and about 0.108 in the report's baseline, before the bequests that the docstrings credit with the whole residual (0.02--0.04, with ``NFA$/Y$ near $-5$'', which the code cannot produce since $K/Y = 3.667$ bounds $(A - K)/Y$ below at $-3.667$). In the results checker the $r\,\mathrm{NFA}$ term is always zero because \code{params['r']} is never written by \code{run_fiscal_figures.py} nor by the checker's configuration merge, so the checker always tests the closed form. Consequence: the one FAIL-tier identity check added at \code{3784ae5} fails a run on an identity the model does not satisfy; its threshold of 0.08 would trip for reasons unrelated to a normalisation error. Verified by reading both implementations; the algebra is agent B3's derivation, re-read. The CPU run at $n_{sim} = 200$ (Appendix~C) measures the open form at $-0.197$ of output at $t = 0$ in the report's baseline, of which the identity's known terms account for $-0.162$, leaving $-0.035$ for the unmeasured accidental bequests.

\textbf{M4. No rest point for debt; the terminal check ignores $B$; the $\tau^l$ target is the baseline's drifted ratio.} With $r_B = 0.019 > \Gamma_T - 1 = 0.017$ the recursion \eqref{eq:debt} has root 1.00197 on the terminal path and a stationary $B/Y$ needs a primary surplus of 0.197\% of output per unit of $B/Y$; during the transition $\Gamma_t - 1 - r_B$ runs from $-0.0069$ to $-0.002$. \code{_check_terminal_convergence} tracks $K$, $K^g$, $L$, $Y$, $C$, $A$, NFA and $S$ and not $B$ (\code{fiscal_experiments.py:411--442}), so ``terminal converged'' is compatible with a drifting debt ratio; \code{run_fiscal_figures.py:326--328} sets the tax-financed target to the baseline's own terminal $B/Y$. Consequence: the reported $\Delta\tau^l$ is defined relative to wherever the baseline's debt ratio has drifted to at $t = 180$ and depends on the horizon. The last logged baseline (13:41, on the pre-edit code) shows pensions$/Y$ rising from 0.157 to 0.289 at $t = 25$ and settling at 0.199, revenue from 0.328 to 0.357; adding the fixed ratios gives a terminal primary deficit of about 0.3\% of output (back-of-envelope). The CPU run at $n_{sim} = 200$ gives a terminal full primary deficit of about 0.06\% of output (Appendix~C). Verified by reading; the baseline $B/Y$ path on the current code does not exist.

\textbf{M5. The bequest circuit is open (July item 2, unchanged).} \code{recompute_bequests} defaults to \code{False} (\code{fiscal_experiments.py:130}) and none of the seven scenarios in \code{run_fiscal_figures.py} sets it; the redistribution loop runs only when it is set (\code{olg_transition.py:2064--2067}); the transfer to the entering cohort is zero without it (\code{:1181--1185}); \code{bequest_transfers} is in the budget dictionary but in neither revenue nor spending (\code{:1784}, \code{:1824--1825}); $\tau^{beq} = 0$; the calibration never sets a transfer. The wealth of the dead (about 3.8\% of output per period in the July runs; not re-measured) leaves the economy, and $\beta$ was fitted with that outflow. The recorded bequest is the dying household's beginning-of-period $a$ rather than the $(1+g)a'$ it carried out of the period (\code{lifecycle_perfect_foresight.py:1262}; JAX \code{:802}), and when the loop is on the transfer goes to the entering cohort of the same period with no $\Gamma_t$ scaling and never to cohorts born before $t = 0$. The draft states that bequests are taxed and redistributed. Verified. Closed on 2026-10-02 by taxing bequests away in full and correcting the measure (\S\ref{sec:postaudit}).

\textbf{M6. Two scripts fail on their own on the sixth parameter.} \code{normalize_A_tfp.py:93} and \code{pin_baseline_closure.py:60} build $\theta$ as \code{[theta_dict[p.name] for p in spec.params]}; the stored $\theta$ lacks \code{ui_replacement_rate}, so both raise \code{KeyError} before any solve. \code{theta_from_config} exists for this case and is used by \code{fill_report.py}, \code{diag_ss_vs_transition.py} and \code{diag_bequest_decomp.py}. Inside \code{run_scale_loop.sh} the failure does not occur, because \code{calibrate.py} rewrites $\theta$ with all six keys on convergence before the two scripts run; re-setting the balancing item on its own, which is what the floor change calls for, fails. Verified by reading both lines and the configuration.

\textbf{M7. The scale loop never evaluates the written $(\theta, A)$ pair.} Described in \S2.3: the convergence test uses the residual at the pre-update $A$, after which $A$ is still rewritten; on 2026-10-01 the pass margin was $0.4\times 10^{-3}$ and the moments at the written pair deviate by up to 1.1\% ($A/Y$). Verified in \code{run_scale_loop.sh:55--61} and the log.

\textbf{M8. The report tables are internally inconsistent.} (a) The implied-statistics rows ``contribution base$/Y$'' and ``contribution rate'' are the identities $1 - \alpha - \mathrm{UI}/Y$ and $\tau^p$ (\code{fill_report.py:261--273}), presented against data 0.570 and 0.228. (b) UI$/Y$ is filed under ``Not targeted'' (\code{UNTARGETED = ['ui_over_Y']}, \code{:121}) while the configuration targets it, and on the markdown path its ``\% dev'' cell prints the absolute gap. (c) The parameter table prints $n = 0.00$ from \code{external_params.pop_growth} (\code{:60}) and the header of the same run prints $n = -0.0048$ (\code{:533}); $\Gamma - 1 = 0.01700$ in the header and the growth table's Level column use $n = 0$, which equals $\Gamma_T - 1$ only because the configuration's terminal rate coincides with the demography file's $n_\infty$. (d) $b_{\min} = 0.17$ is printed beside $\theta$ fitted at 0.15, and five SMM rows beside six configured parameters. Verified in the generator and the bodies.

\textbf{M9. $B_0$ and the $I_g$ shock are sized off a 50-household preliminary run (July item 4, unchanged).} \code{run_fiscal_figures.py:102--106, :123} and, for the shock, \code{:266--267}; the full runs use 2\,000 households per cohort. $B_0/Y_0$ differs from 1.6428 by the preliminary run's sampling error, which propagates through the whole $B/Y$ path and the tax-financed target. Verified.

\subsubsection*{Minor}

\begin{enumerate}[label=m\arabic*.]
\item \textbf{Dating of the tax-financed target.} \code{run_fiscal_figures.py:326} sets the target as $B_{T_{bal}}/Y_{T_{bal}-1}$ with a comment saying it matches the residual's convention; \code{_balance_residual} now evaluates $B_{T_{bal}-1}/Y_{T_{bal}-1}$ (\code{fiscal_experiments.py:335--337}, changed 2026-10-01). The counterfactual is required to reach at $t = 179$ the ratio the baseline reaches at $t = 180$; bias in $\Delta\tau^l$ of order 0.01 pp. Verified.
\item \textbf{G multiplier when G is a share of output.} $G_t = 0.13\,Y_t$ in the baseline and $0.15\,Y^{cf}_t$ in the shock, so $\Delta G_t = 0.02 Y^{cf}_t + 0.13\,\Delta Y_t$ and the reported multiplier is $m/(1 + 0.13 m)$ for a true impact $m$; the $I_g$ shock is a level. Read, not re-derived.
\item \textbf{The balancing-item script writes the configuration-ratio value.} \code{pin_baseline_closure.py:76--95} recomputes $I_g/Y$ from the level and prints it, then writes \code{ratios['closure_other_over_Y']}, formed from \code{fiscal.I_g_over_Y} (\code{:104}); identical to five decimals today. Verified.
\item \textbf{$z_{last}$ convention (July item 3, unchanged).} UI lasts one period per spell, entrants drawn unemployed get none, and a worker unemployed at $J_R$ retires on the $(1-\lambda)$ term; four sites, both implementations. The draft now describes the one-period UI; the retirement consequence is not in the draft.
\item \textbf{Silent defaults (July item 6, unchanged).} $\lambda$ from a constant $\hat\rho = 0.95$ (\code{calibrate.py:1342}), not from \code{edu_params.rho_y}; $\tau^{beq}$ from the dataclass default; neither key in the configuration.
\item \textbf{Pre-2023 wage path is \code{None} in every fiscal run.} The baseline paths that pre-2023 cohorts keep in a counterfactual are assembled before \code{base_paths['w_path']} is set and on a local copy (\code{fiscal_experiments.py:1226, :1244--1248, :1278}); harmless because the baseline decision rules are stored before any counterfactual runs and $w_0$ is the same in every scenario, and each scenario call re-runs the baseline (cost only).
\item \textbf{Pension fund at $r_B$.} The memorandum recursion accrues at $r_B$ (\code{olg_transition.py:2305--2317}); the draft and the plan write $r_t$. Feeds nothing.
\item \textbf{Implementation and seed details.} The JAX single-model \code{simulate} ignores \code{T_sim}, so the cohort-by-cohort construction simulates every cohort over the full horizon on JAX and reads row $j$ (no numerical effect); the SMM seeds cohort $j$ with $42 + j$ for all three education types (common random numbers across education); the JAX simulation of cohorts solved together and the single-model JAX simulation split keys differently for $n_\alpha > 1$; the single-model JAX solve drops \code{bequest_lumpsum} while the kernel for cohorts solved together carries it (inert at zero); dead households' $\ell$ is 1 in NumPy and 0 in JAX (not aggregated).
\item \textbf{Records.} The report prints $n_{sim} = 10\,000$ where the cohort-by-cohort construction uses 2\,000 per cohort; \code{theta_metadata.source_report} names a directory that holds only July reports; the normalisation and balancing-item docstrings say ``one lifecycle problem''; \code{fill_report.py:119} says the interest ``data'' is 0.0312 while the configuration holds 0.0339; the configuration's survival vector is the 2020 table; the docstrings' ``0.47 spread in the probability of reaching 84'' is 0.25 on the demography file (0.505--0.759).
\item \textbf{Test fixtures with a positive consumption floor.} The constructor now refuses \code{transfer_floor > 0}, so the two household-isomorphism tests fail at fixture construction and three further tests that build such fixtures (\code{test_olg_transition.py:1160, :1410, :1436}) would too; the isomorphism itself holds exactly with the floor at zero (Appendix~C).
\item \textbf{Entry points that cannot complete.} \code{python olg_transition.py} without flags builds $n_h = 2$ and \code{_alive_fraction} refuses it; the \code{lifecycle_perfect_foresight.py --test} block unpacks 19 names from a 22-tuple. Neither is a production path (Appendix~B).
\item \textbf{Two live configurations diverge.} \code{calibration_input_GR_rB0.json} (July: $K^g = 0.745$, 60 periods, no demography file) is tracked, read by nothing in \code{code/}, and still describes the July mechanism of item 5.
\item \textbf{Hours of the unemployed in the simulated sample.} The sample records $\ell = 1$ for unemployed households (a placeholder from the decision arrays) and $0$ for retirees; the hours moment excludes the unemployed and $L_t$ is built from income, so the value reaches no aggregate or moment.
\end{enumerate}

\subsection{The six items of the July open-issues list at HEAD}

\begin{center}\footnotesize
\begin{tabular}{@{}clp{9.6cm}@{}}
\toprule
Item & Status & Evidence at HEAD \\
\midrule
1 $L$ includes UI & still present & M2; four sites, both implementations; only line numbers moved. \\
2 bequest circuit open & still present & M5; seven (not six) scenarios, none sets the flag; the configuration path of the command line also defaults to off. \\
3 $z_{last}$ = previous state incl.\ 0 & still present & m4; four sites, both implementations. \\
4 $B_0$ from the preliminary run & still present & M9; the $I_g$ shock level shares it. \\
5 $Y_{ss} = 1$ not carried to $t = 0$ & changed & The mechanism the note names (configuration survival vector against cohort schedules) is gone for the production configuration: both constructions now take cohort schedules from the demography file and the calibration weights the measured cross-section. The last production-scale measurement (13:41, single-lifecycle, $n_{sim} = 2\,000$) was $-9.5\%$; the CPU run on the cohort-by-cohort construction at $n_{sim} = 200$ gives $+1.2\%$ (Appendix~C), one draw. The aggregation conventions, seeds and random-number generators still differ. For the rB0 configuration the item stands as written. $K^g$ is now 0.703, not 0.745. \\
6 silent defaults & still present & m5; the note's ``derives $\lambda$ from $\rho_z$'' overstates it: a constant 0.95 that coincides with $\rho_z$. \\
\bottomrule
\end{tabular}
\end{center}

\subsection{Code against the draft's model section}

The draft (\code{docs/DSA-LSA model.tex} at submodule commit \code{da44637}, 2026-09-21) predates the branch. Disagreements that \code{code/docs/TREND_GROWTH_PLAN.md} settles are booked as settled changes: CRRA utility with curvature $\sigma$ (code: $\log c$); population growing at a constant $g_N$ with weights $\Lambda^{-j}$ (code: measured 2023 cross-section and the EUROPOP path); $a'$ without the growth factor (code: $(1+g)a'$); the public-capital and debt laws without $\Gamma$ (code: divided by $\Gamma_t$); no statement on who holds $B$ (plan: external official-sector debt at $r_B$). Every other disagreement is a finding:

\begin{itemize}
\item The draft's retired Bellman equation continues with $V^R_{j+1}(z_{last}, a')$; the code continues with $z_{last} = 0$ (C1).
\item The draft's $L_t = \int \kappa_j z e^{\alpha}\ell\, d\mu$; the code's $L_t$ includes UI$/w$ (M2).
\item The draft taxes bequests at $\tau^{beq}$ and redistributes them to newborns; the code records them and lets them leave (M5).
\item The draft enforces ``a minimum consumption floor via a lump-sum transfer when needed''; the code has no floor and refuses a positive one (m10). The default rule of M1 is the only thing that happens at the bottom.
\item The draft's pension fund accrues at $r_t$; the code's at $r_B$ (m7).
\item The draft describes UI as ``zero once a spell runs beyond one period'', which matches the code; it does not state that the $\lambda$ term of the pension is zero for a household unemployed at $J_R$ (m4).
\end{itemize}
The pension formula, the four tax bases, the health line, the production function and both firm conditions, $\mathrm{NFA} = A - K - B$, the primary-deficit definition with $D_t$ and $O_t$, and the cohort-specific survival along calendar years agree between draft and code.

\subsection{Dead-code inventory}

Agent B6 parsed the 28 Python files with \code{ast} (611 definitions: 324 in non-test modules, 287 collected by pytest), indexed every name and attribute reference with its file and line, kept string-only hits apart, counted the four shell scripts as callers, and confirmed each candidate by reading its call sites against the production configuration and scripts. Counts of tabulated rows (a row can cover several definitions): (a) unreachable 6; (b) test-only 11 rows, 15 definitions; (c) switched off in production 31; (d) dead inside live code 39; (e) configuration keys without production effect 13; (f) scripts no shell script calls whose outputs nothing reads 5; plus 3 observations that are not dead code. The repeated search reproduced the five unreachable definitions (\code{solve_labor_hours_jax}, \code{_solve_labor_hours}, \code{_compute_wage_path_njit}, \code{_get_cached_cohort_panel}, \code{compute_aggregates}: definition only, one docstring mention); the sixth item is a set of branches of \code{compare_scenarios} that no caller reaches. Items in (d) with a bearing on correctness rather than tidiness: a stored simulation result (\code{_birth_sim_cache}) written and never read under JAX; the budget keys \code{debt_service}, \code{new_borrowing}, \code{fiscal_deficit} identically zero or duplicate in production yet written to the results JSON; \code{simulation.n_sim} printed but not used on the cohort-by-cohort construction; the pension-weight constant of m5; the duplicated aggregate implementations (\code{_compute_ss_aggregates} against \code{compute_fiscal_ratios}; the regeneration script against the figure script). Whether any item is deleted is not part of this report. The full inventory, with the pattern searched and the result per item, is Appendix~B; the raw 611-row table is \code{audit_2026-10-01/dead_code_inventory.csv}.

\subsection{Not verified}

\begin{itemize}
\item The production-scale magnitude of C1 and M1 (requires the 60-cohort solve at the production grid; the agents' checks were small models).
\item The $t = 0$ output gap, the goods-market residual and the flatness column at the production $n_{sim} = 2\,000$: C1 of the plan was not run (no GPU in the session). A CPU run at $n_{sim} = 200$ measured them once (Appendix~C): gap $+1.2\%$, open residual $-0.197$, trends within $0.07\%$ per year; not repeated across seeds.
\item Predetermination of initial wealth across scenarios under the Step 0 demography at production scale (C1b): not run. The mechanism passes on the small test economy on both implementations.
\item Numerical agreement of the two implementations at production scale; the construction and implementation behind the recorded $\theta$ (inferred, not logged); the baseline $B/Y$ path and terminal primary balance on the current code; the size of the accidental-bequest outflow and of the sampling error in $Y_0$.
\item Whether the user-level skills outside \code{code/} call \code{eval_fiscal_results.py} and \code{regen_fiscal_figures_from_json.py} (B6 classified them as manual command-line tools).
\end{itemize}

\subsection{Comments received on the first draft (2026-10-02)}\label{sec:comments}

Comments that overlap a finding above were folded into it: bequests must be taxed away or transferred (M5); labour input should be the sum of hours in efficiency units, not of income (M2, where dividing wage income by the wage gives exactly that, so the benefit term is the whole discrepancy); unemployment benefits do not belong in the resource constraint (M2 and M3). Three do not overlap and are recorded here with the facts the code gives.

\begin{enumerate}[label=c\arabic*.]
\item \textbf{Should $B$ be counted in NFA?} $\mathrm{NFA}_t = A_t - K_t - B_t$ is the economy's net foreign asset position: household foreign assets $A_t - K_t$ less the government's external debt. Under the assumption that $B$ is held abroad it must be netted out, with a negative sign, to obtain the country's position; the household sector's own position is $A_t - K_t$. The code uses the economy-wide definition for the NFA target of the tax-financed experiments and for the current account, and the household definition nowhere; the two are consistent and the report labels which is which.
\item \textbf{How was the job-finding probability calibrated?} $f = 0.5$ is a constant in the configuration (\code{external_params.job_finding_rate}) with no recorded source in the code, the configuration or the calibration reports; the separation rates are derived from it and the education unemployment rates. Whether $f$ was ever calibrated, and to what, is not recorded.
\item \textbf{Hours should not be restricted to $[0, 1]$.} The solver restricts the labour first-order condition's solution to $[0, 1]$ (\code{lifecycle_perfect_foresight.py:841--891}; \code{lifecycle_jax.py:56--93}), so where either bound binds the condition holds as an inequality. Whether the upper bound binds for any household-period at the production parameters was not measured; average hours among the employed are 0.41.
\end{enumerate}

\subsection{Changes made after the audit (2026-10-02)}\label{sec:postaudit}

Made at the user's request after the review of the first draft; they postdate the audited commit. The bequest change is committed as \code{ff76dd4}; the rest is in the working tree at the time of writing. Decisions taken by the user: a booked consumption floor for the non-positive-resource states (M1); debt tracked and reported rather than pinned (M4); no upper bound on hours (c3); the job-finding probability calibrated to unemployment duration (c2).

\begin{enumerate}[label=p\arabic*.]
\item \textbf{Accidental bequests taxed away in full} (M5). \code{external_params.tau_beq = 1.0}; the base-year budget in \code{calibrate.py} now carries the bequest line, so the balancing item and the transition see the same revenue. The recorded bequest is the wealth the dying household carried out of the period, $(1+g)a'$, in both implementations. Checked as reported in the first version of this section; committed as \code{ff76dd4}.
\item \textbf{C1, retired continuation value.} Both solvers read $V_{j+1}$ at the retiree's own $z_{last}$ (\code{lifecycle_perfect_foresight.py}, retired branch of \code{_solve_state_choice}; \code{lifecycle_jax.py}, \code{solve_period_jax}, where the $y_{last}$ axis is kept through the health expectation). Check: on a small model the retired value varies with $z_{last}$, is nondecreasing in it, and the two implementations agree to $10^{-14}$ with identical savings rules.
\item \textbf{M2, labour input.} $L = (\text{wage income} + \mathrm{UI} - \mathrm{UI})/w$ in the transition, in the unused \code{compute_aggregates}, and in both aggregators of \code{calibrate.py}. Check: $wL$ equals the wage bill, payroll tax over its rate, to $10^{-10}$ in a transition with a positive payroll rate.
\item \textbf{C2, horizon beyond the entering-cohort table.} \code{_entrant_weights} extrapolates the entering cohort past the table at the terminal rate $n_\infty$, which is exact because the table ends after the ramp; \code{growth_factors} already clipped. Check: period 199 returns the same weights as 187 on the production configuration.
\item \textbf{M6 and m3, the two standalone scripts.} \code{normalize_A_tfp.py} and \code{pin_baseline_closure.py} read $\theta$ with \code{theta_from_config}; the closure script writes the value built on the level-based $I_g/Y$.
\item \textbf{M1, consumption floor, booked.} \code{external_params.transfer_floor = 0.0807}: the Greek guaranteed minimum income for a single adult, EUR 200 a month, over output per person aged 25--84 (\code{build_transfer_floor_GR.py}, \code{data/transfer_floor_GR.json}; the amount is to be checked against the base-year vintage). The budget function returns the top-up it granted; both simulations record it per household (\code{transfer_sim}, the 23rd panel array), the per-age means carry it as a 12th element, the transition's budget books it as an outlay (\code{transfers}) and the base-year budget nets it in the primary balance. The two refusals of a positive floor are gone; the JAX wrapper refuses a floor together with schooling, whose child costs the simulated transfer does not replicate. Check: with a floor of 0.08 on a small model the household budget identity holds to $10^{-14}$ on both implementations with the transfer included, and the transition's total spending equals the sum of its lines including transfers. With the floor above the largest out-of-pocket medical cost, the non-positive-resource states of M1 no longer occur.
\item \textbf{M4, debt in the terminal check.} \code{_check_terminal_convergence} takes the debt path and reports its drift at the slow tolerance; the three fiscal runners pass it. No rest-point rule is imposed, by decision.
\item \textbf{M7, the scale loop.} \code{run_scale_loop.sh} converges only if, in addition to the pre-update residual, every targeted moment at the written $(\theta, A)$ pair is within \code{MOM_TOL} (default 0.5\%), read from the normalisation's own printout.
\item \textbf{M9 and m1, the fiscal driver.} The preliminary run is at the full sample size, so $B_0/Y_0$ equals the configured ratio exactly and the $I_g$ shock is sized off the same output; the tax-financed target is dated $B_{T_{bal}-1}/Y_{T_{bal}-1}$ as the residual is; the results now stamp $r$, $\kappa$, $\tau^{beq}$ and the floor.
\item \textbf{M3, the resource-constraint checks.} The report generator's baseline receives the spending shares and computes the budget; its residual is $C - (Y + r\,\mathrm{NFA}_p - I - G - I^g - D - O - M - \Delta\mathrm{NFA}_p + \mathrm{PD})$ on the budget's own lines, with $M$ total medical spending, warning above 1\%. The results checker's residual is $C - (Y + r\,\mathrm{NFA} + (r - r_B)B - I - \text{purchases} - M - \Delta\mathrm{NFA})$, with $r$ now read from the results and the debt path passed in, failing above 2\%. Neither has been run on a production transition yet.
\item \textbf{M8, the report tables.} The parameter table prints $n_0$ and $n_\infty$ as separate rows, every parameter the configuration lists for the SMM (a never-fitted one shows no value), the floor, $\tau^{beq}$ and $f$; UI$/Y$ is not filed as untargeted when it is a target, and the markdown route recomputes the relative deviation; the two identity rows of the implied table are labelled as such; the header carries $\Gamma_0 - 1$ and $\Gamma_T - 1$ and the growth table's level column uses $\Gamma_T - 1$.
\item \textbf{c3, hours.} Both labour solvers bracket upward by doubling until the condition's residual turns positive, so hours above one are admitted and the condition holds with equality at every interior solution. Check: with a small disutility weight both implementations choose hours above one, and a direct solve returns a residual of zero at $\ell = 3.84$.
\item \textbf{c2, job-finding probability.} \code{build_job_finding_GR.py} derives $f = 1 - \text{long-term share of unemployment}$ from Eurostat \code{une_ltu_a} (ages 15--74, both sexes; \code{data/job_finding_GR.json}); the user chose the base year: share 0.560 in 2023, $f = 0.440$, written to the configuration (it was 0.5, unsourced). Alternatives recorded in the file: 0.385 over 2015--2024, 0.406 over 2019--2024, 0.464 for 2024.
\item \textbf{m13, hours of the unemployed.} The simulated panel records $\ell = 0$ for the unemployed on both implementations; the two labour-supply tests that asserted the old placeholder were updated.
\item \textbf{m10, test fixtures.} With the floor booked the constructor accepts a positive floor again, so the household-isomorphism tests run; six regression tests for the items above were added (\code{TestAuditFixes20261002}).
\item \textbf{Found by the full test suite after the changes, and fixed.} The weighted Gini coefficient (\code{calibrate.compute_gini} with weights) summed the cumulative income share after each observation instead of the trapezoid mean of the shares before and after it, so every weighted Gini was understated (0.067 instead of 0.267 for the values 1 to 5 with equal weights; negative on skewed weights). Pre-existing: identical on the audited commit, and the suite's own test asserted agreement only to 0.05. The wealth, earnings, income and disposable-income Gini rows of every calibration report to date are affected; none is a target. Fixed to the trapezoid rule, with an exact test against an expanded sample. A second failing test, which asserted that a higher pension floor raises retiree consumption, encoded a claim the model does not imply: with the continuation value now read at the household's own last income state, the fixture smooths by consuming more before retirement and 0.05\% less after it. The test now asserts what the theory does imply, that the value function weakly rises at every state and the floor binds.
\item \textbf{Left open.} Everything above changes what the SMM fits to, so $\theta$, $A$ and the balancing item must be re-fitted (C3); the resource-constraint checks and the fiscal pipeline have not yet been exercised on a production transition; the \code{retirement_window} path, the HSV schedule and the other switched-off features of Appendix~B were not touched.
\end{enumerate}

%====================================================================
\appendix
\section{Code-to-equation correspondence}

Line numbers are from the files at \code{3784ae5}. PF = \code{lifecycle_perfect_foresight.py}, JX = \code{lifecycle_jax.py}, OT = \code{olg_transition.py}, FE = \code{fiscal_experiments.py}, CA = \code{calibrate.py}, RF = \code{run_fiscal_figures.py}, FR = \code{reports/fill_report.py}.

\footnotesize
\begin{longtable}{@{}p{3.6cm}p{5.3cm}p{3.6cm}p{3.6cm}@{}}
\toprule
Object & Equation (Part I) & NumPy / scripts & JAX \\
\midrule
\endhead
Utility, disutility & $\log c$; $\nu \ell^{1+\varphi}/(1+\varphi)$ & PF:629--634, 893--897 & JX:35--41, 96--98 \\
Discounting with survival & $\beta\,\pi_j\,\mathbb{E}V_{j+1}$ & PF:946, 1012 & JX:338--339, 345 \\
Working budget & \eqref{eq:bcw}; $(1+g)a'$ & PF:742--768, 974 & JX:163--190, 288 \\
Retired budget, pension & \eqref{eq:bcr}, \eqref{eq:pens} & PF:724--740, 1194--1200 & JX:140--161, 713--720 \\
UI & $T^{UI}$ in \eqref{eq:pens} & PF:743--747, 1215--1218 & JX:164--168, 723--727 \\
Tax bases & $\tau^c, \tau^l, \tau^p, \tau^k$ & PF:751, 758, 739, 761--763, 1235--1254 & JX:170, 176, 159, 183--185, 743--763 \\
Labour FOC and bounds & $\nu\ell^{\varphi}(1+\tau^c) = \mathrm{MW}/c$ on $[0,1]$ & PF:841--891 & JX:56--93, 298--310 \\
Working continuation & $\sum_{z'} P_z V_{j+1}(z', z, a')$ & PF:997--1002 & JX:320--330 \\
Retired continuation (C1) & $V_{j+1}(0, a')$ & PF:993 & JX:332 \\
Terminal age; default rule (M1) & $a' = 0$; $c = 10^{-10}$, no continuation & PF:953--965; 1020--1024 & JX:366--439; 354--361 \\
Asset grid & $a_i = 200(i/99)^{1.5}$ & PF:435--447 & JX:1032 (same array) \\
Income process & Tauchen + $z = 0$; $f$, $s_e$; $\alpha$ discretisation & PF:496--534, 538--559 & copied, JX:1015, 1032--1037 \\
Entry state & $a = 0$, $z \sim$ stationary, $z_{last} = z$ & PF:1137--1155; OT:630--642 & JX:1201--1225 \\
$z_{last}$ update; freeze & previous state incl.\ 0 & PF:1270--1271 & JX:790 \\
Death, bequest record & $1 - \pi_j$; $a$ recorded & PF:1256--1264 & JX:797--803 \\
Cohort survival schedule & $\pi_j(y)$, clamped years & OT:1023--1040, 1047--1066 & same (schedules stacked, OT:440--444, 689--692) \\
Cohort paths before 2023 & period-0 values & OT:58--80, 1154--1166 & same \\
Baseline rules kept in counterfactuals & pre-2023 ages & OT:1202--1273 & OT:1308--1345 \\
Population weights, living share & \S1.1 & OT:938--956, 973--1021 & same \\
$\Gamma_t$ & $(1+g)(1+n_t)$, clipped at the table end & OT:905--936 & same \\
Aggregates $A, C, L$ & per capita; $L = \Sigma(\cdot)/w$ & OT:806--823, 2163--2173 & same \\
Firm conditions, $K$, $Y$ & \S1.5 & OT:2002--2004, 2185--2198; CA:1069--1096 & same \\
Public capital & $K^g_{t+1} = [(1-\delta_g)K^g_t + I^g_t]/\Gamma_t$ & OT:1972--1978 & same \\
Baseline $I^g$ level & $(\delta_g + \Gamma_t - 1)K^g$ & RF:98--99; FR:540--541 & --- \\
Government budget & \eqref{eq:pd} & OT:1735--1851 & same \\
Debt law & \eqref{eq:debt} & FE:264--297 & --- \\
$B_0$ (M9) & $\bar b\,Y_0$ & RF:102--106, 123 & --- \\
NFA, CA & $A - K - B$; $\Gamma_t \mathrm{NFA}_{t+1} - \mathrm{NFA}_t$ & FE:650--667, 755--758 & --- \\
Pension fund (memorandum) & $[(1+r_B)S + \mathrm{Rev}^p - \mathrm{PENS}]/\Gamma_t$ & OT:2305--2317 & --- \\
Balance residual, targets & $B_{T_{bal}-1}/Y_{T_{bal}-1}$; NFA; rest point & FE:326--357; RF:320--334 & --- \\
Terminal-convergence check & $K, K^g, L, Y, C, A$, NFA, $S$ & FE:406--461 & --- \\
Horizon extension (C2) & 20 periods & RF:218; FE:217--243, 730--736 & --- \\
Goods-market residual (M3) & \S1.7 & FR:302--364; \code{eval_fiscal_results.py:225--262} & --- \\
SMM cross-section (cohort by cohort) & 60 solves per education type & CA:748--817, 826--836 & same classes \\
Moments, weights & among the living, 2023 shares & CA:328--353, 587--695, 1174--1189 & --- \\
Scale loop, balancing item & \S2.3 & \code{run_scale_loop.sh:23--68}; \code{pin_baseline_closure.py:60--120} & --- \\
\bottomrule
\end{longtable}
\normalsize

\section{Dead-code inventory}

Agent B6's inventory, reproduced from its report; classes (a)--(f) as defined in the audit plan. Line numbers are from the files at \code{3784ae5}. No view on deletion.

\footnotesize
\subsection*{Method}
\input{\tabdir/dead_code_method.tex}
\subsection*{Items}
\input{\tabdir/dead_code_tables.tex}
\subsection*{Not settled by reading}
\input{\tabdir/dead_code_unsettled.tex}
\normalsize

\section{Checks run and their outcomes}

All local, on an 8-core Apple CPU with 24 GB, JAX 0.8.0 on CPU in float64, the \code{jax-arm} virtual environment, on 2026-10-01 between 22:50 and 23:30 local time, plus the overnight run described below. Scripts and logs are in the session scratchpad, not in the repository.

\begin{longtable}{@{}p{4.2cm}p{11.3cm}@{}}
\toprule
Check & Outcome \\
\midrule
\endhead
C0, cost of the cohort-by-cohort construction & One education group, production grid ($J = 60$, $n_a = 100$, $n_y = 5$, $n_\alpha = 5$), JAX, $n_{sim} = 200$: \code{base_year_cross_section} 956 s, 15.9 s per cohort; a single solve 11.5 s (no compile reuse across instances), ratio about 83. Scaling: one SMM evaluation on this construction is three groups, about 48 min on this CPU; the A100 did a three-group single-lifecycle evaluation in about 1.3 s against 35 s here (about 30 times faster), so one cohort-by-cohort cross-section is about 3 min of A100 time, one SMM evaluation about 3 min, and the 169-iteration round of 12:31 would be of the order of 10 hours; C2 is tens of GPU hours. C1 is about 10 min of cross-sections plus two transitions, which the 13:41 chain completed within about an hour. \\
Spot check 1, living share & \code{_alive_fraction(t)} equals the demography file's living over ever-entered to $2\times 10^{-16}$ at $t \in \{0, 1, 10, 30, 60, 100, 150, 179\}$; the weights times cumulative survival sum to 1 at each (true by construction, as the test file notes). \ok \\
Spot check 2, $\Gamma_t$ & \code{growth_factors(180)[t]} $= (1+g)\,N_{t+1}/N_t$ exactly (forward convention); $\Gamma_t = 1.017$ for all $t \ge 156$; the model's population from entering cohorts and the code's cohort survival equals the demography file's population to $2\times 10^{-16}$ at six dates. \ok \\
Spot check 3, survival schedules & \code{_cohort_survival_schedule} equals the demography file's $p(\text{year}, j)$ exactly for cohorts born 1964, 1993, 2023, 2073, 2143, and the historical table where it covers; the table is constant after 2100 and equals the EUROPOP 2100 row. Probability of reaching 84: 0.514, 0.633, 0.759, 0.840, 0.842. \ok \\
Spot check 4, the two constructions agree under one schedule & Covered by \code{TestBaseYearCrossSection::test_shared_schedule_reproduces_the_single_solve_exactly}, passed. \ok \\
pytest selection (26 tests, serial) & 23 passed, 1 skipped (cost ratio of cohorts solved together, GPU only), 2 failed: the two household-isomorphism cases fail at fixture construction because the fixture sets \code{transfer_floor = 0.05} and the constructor now raises; stale fixture, not a model failure. 320 s, of which 296 s in \code{test_kg_flat_at_stationary_investment} (NumPy transition). \\
Isomorphism with the floor at zero & Same fixture otherwise, both implementations: \code{a_policy} identical; $\max|\Delta c| \le 8.9\times 10^{-15}$, $\max|\Delta \ell| = 2.2\times 10^{-16}$, $\max|\Delta V| \le 5.3\times 10^{-15}$. The growth problem on grid $G$ at $r$ equals the no-growth problem on $(1+g)G$ at $\tilde{r} = 0.017801$. \ok \\
\code{check_a0_predetermination.py} & Small test economy ($J = 20$, $n_a = 30$, $n_{sim} = 50$, $T_{tr} = 10$, $n_\alpha = 3$), $\tau^l$ and $I_g$ shocks, both implementations: $|A_0^{cf} - A_0^{base}| = 0$ in all four cases; values identical to the 13:41 record. Says nothing about the production economy. \ok \\
Entering-cohort table bound (C2) & \code{_entrant_weights(187)} returns; \code{(188)} and \code{(199)} raise; \code{growth_factors(200)} clips. \verified \\
Dead-code search repeat & Five unreachable definitions reproduced (definition only); \code{compare_scenarios} is called with explicit variables at every site. \verified \\
C1 (GPU) & Not run: no instance in the session. \\
C1b (GPU) & Not run. \\
C2 (GPU) & Not run; sized by C0 at tens of GPU hours. \\
Reduced-scale substitute for C1 & Cohort-by-cohort cross-section and one baseline transition at $n_{sim} = 200$, JAX on CPU, production configuration and $\theta$ (\code{c1_local.py}): completed, 2\,866 s + 10\,969 s; results in the note below the table. \\
\code{check_a0_predetermination.py} after the 2026-10-02 change & Re-run after the bequest edits of \S\ref{sec:postaudit}: all four cases $|A_0^{cf} - A_0^{base}| = 0$, values identical to the pre-change run to ten decimals (the harness has $\tau^{beq} = 0$, and the measure does not enter wealth); the vmapped simulation with the new positional argument runs on both implementations. \ok \\
\bottomrule
\end{longtable}

\input{\tabdir/c1_local_status.tex}

\end{document}
